\documentclass[11pt]{article}
\usepackage{amsmath,amssymb,amsthm,geometry}
\geometry{margin=1in}
\newtheorem{theorem}{Theorem}
\newtheorem{lemma}[theorem]{Lemma}
\newcommand{\Q}{\mathbb Q}\newcommand{\Z}{\mathbb Z}
\DeclareMathOperator{\Tr}{Tr}
\title{An affine fourth label in the reciprocal quartic construction}
\author{Research note; AI generated}
\date{14 November 2026}
\begin{document}\maketitle

This note supplies the arithmetic argument for replacing the fourth
seven-field root block of Brox's 204-coordinate family by a five-field
reciprocal pair. It is relative to precisely the scalar, pole,
contact-height, elliptic group, and fixed-data prime-selection premises
retained in \texttt{quartic204\_20261114\_v27.tex}, at repository commit
\texttt{d337a5500e2531139445bd7e30f966f62c4db03a}.
It does not independently certify those premises. The quadratic-table
construction and its independent verification are separate obligations.
All counts include the rational input.

Put $K=\Q(\rho)$, $\rho^2=2$, $d=3-\rho$, $D=d^2$, and
$E:y^2=x^3-Dx$. Retain $P=2R_0$ and $2E(K)=\Z P$.
For $Q=[m]R_0+T_0$, $T_0\in E[2](K)$, put $\nu(Q)=m/2$.
The first and third pairs, scalar block, and reference point are unchanged.
In particular $Q_2=[a]P$, where $a\in\Z\setminus\{0\}$, and
\[
 \beta=51X^2+2504,\quad R>100,\quad
 c=|a|\ge149^{L_N},\quad
 h(\beta)\le C_3 149^{-L_N}c(1+\log c/3),
 \qquad L_N=\left\lfloor\frac{9(2N-1)}4\right\rfloor+1.
\]
Here $h$ is logarithmic Weil height and $C_3$ is the fixed source constant.

\section*{The new point block and the affine label}
Use five field variables $f,h,g,k,s$ for the fourth pair, and reuse its
reciprocal-square slots $F,G$. Impose
\[
 F=f^2,\quad G=g^2,\quad h^2=f-DfF,\quad
 s(f-g)=h+k,\quad s(h-k)=1-D(F+fg+G).
\]
The complete reciprocal-pair lemma in the source identifies points
$U_4,V_4\in E(K)\setminus\{(0,0)\}$ with $U_4\ne V_4$ and finite difference
$Q_4=U_4-V_4$, including every tangent and numerator-identity case.
Its recovered coordinates are
\[
 x_4=s^2+D(f+g),\qquad y_4=Dh-s(Df+x_4).
\]
Let $S=y(H)/x(H)$ and $T_H=y(H-P)/x(H-P)$ be the unchanged nonzero
reference parameters. Store the affine ratios
\[
 L_4=4z(U_4)/z(H)+1,\qquad L_4^0=4z(V_4)/z(H-P)+1,
 \qquad z=-x/y,
\]
using the four field quadrics
\[
 h(L_4-1)+4fS=0,\quad (1-DF)(L_4-1)+4hS=0,
\]
\[
 k(L_4^0-1)+4gT_H=0,\quad (1-DG)(L_4^0-1)+4kT_H=0.
\]
The second equation in either pair fixes the correct value $1$ at an
identity numerator. The original nonidentity-two-torsion exclusion is
unchanged: $h=0,f\ne0$ in the first equation would force $S=0$.
Set $L_4'=\sigma L_4^0$ and $C_4=L_4-L_4'$ as before.
The fourth integer label is now
\[
                  b=4\nu(Q_4)+1\in2\Z+1.
\]
Thus $b\ne0$ without a guard, even when $Q_4$ is finite torsion.
In particular it is legitimate that the new chart does not force its
index to be nonzero.

Replace the Lucas expression by
\[
 G_{\rm Luc}=L_3^2-(L_2+1)(L_2-1)L_3L_4+(L_2-1)^2L_4^2-1.
\]
It uses the same two field multiplication outputs as before, now storing
$(L_2-1)L_4$ and $(L_2+1)L_3$. The affine constants do not raise the
quadratic degree. All local factor equations, including the reciprocal
abscissa tests and their exponents, are retained.

Replace the trace aggregate by
\[
 T=\frac14\Tr\bigl(x(Q_1)+x(Q_3)+x(Q_4)+R^2x(Q_2)\bigr),
 \qquad t_0=\lambda,\quad t_1=\lambda+T.
\]
The recovered $x(Q_1),x(Q_3),x(Q_4)$ are quadratic expressions. The old
scalar root expression already stores $R^2x(Q_2)$ quadratically.
At either real embedding a finite elliptic point has $x\ge-d$, whereas
the nonzero doubled scalar point has $x(Q_2)>d$. Consequently
$4T>6R^2-18>0$.

\section*{Height and local conclusions}
Put
\[
 n=\nu(Q_1)\in\tfrac12\Z,\qquad t=2\nu(Q_3)\in\Z,
 \qquad B=\max(1,|n|,|a|,|t|,|b|).
\]
The source weighted-trace theorem applies with weights $1,1,1,R^2$.
Since $b=4\nu(Q_4)+1$, changing the fixed constants gives
\[
                 B^2\le A_0(h(R)+h(T))+J_0.
\]
The affine fourth ratio differs from four times the old physical ratio
only by $1$. Its comparison difference is therefore four times the
unshifted comparison. At every consumed prime the factor $4$ is a unit;
the unchanged local tests and formal-group argument imply
\[
 L_1\equiv n\pmod p,\quad L_2\equiv a\pmod{p^3},\quad
 L_3\equiv t\pmod{p^3},\quad L_4\equiv b\pmod{p^3}.
\]
This statement includes numerator identities and torsion differences,
by the source argument after multiplying the group relations by $2$.
Thus, writing $M_j=M(t_j)$ for the source contact radicals,
\[
 M_j^3\mid\Delta,\qquad M_j\mid J_X\quad(J_X\ne0),\qquad j=0,1,
\]
where
\[
 \Delta=t^2-(a+1)(a-1)tb+(a-1)^2b^2-1\in\Z,
\]
and $J_X$ is the positive absolute reduced numerator of $\beta-n$.
For $B\ge1$, $|\Delta|\le14B^4$ and, when nonzero,
$J_X\le4B\exp h(\beta)$.

The retained scalar normalization still yields
$h(R)\le4h(\beta)+C_RN+C_0$. Hence, exactly as in the source, if
$M_j^k\mid J\ne0$ for both $j$, then
\[
 B^2\le C|J|^{7/(6k)}+C'h(\beta)+C''N+C'''.
\]
All constants are fixed before $N$, $X$, and every witness. The scalar
bound gives $h(\beta)=O(B(1+\log B))$ uniformly in $N$; its forced
lower bound gives $N=O(\log B)$.

\section*{Uniform soundness}
If $\Delta\ne0$, the preceding bound with $k=3$ gives
\[
 B^2\le C_1 B^{14/9}+C_2 B(1+\log B)+C_3'\log B+C_4.
\]
Its right side is $o(B^2)$, so $B$ is bounded by a fixed constant.
For all sufficiently large fixed $N$, the lower bound
$B\ge c\ge149^{L_N}$ excludes this case. Thus $\Delta=0$.

Write $A=a+1$, $c'=|A|$, and
$w=\operatorname{sgn}(A)(A-2)b$. For sufficiently large $N$, $A$ and
$A-2$ are nonzero. The odd integer $b$ is nonzero, so $w\ne0$, and
\[
                     t^2-c'tw+w^2=1.
\]
The source elementary norm-one-unit argument, with recurrence
$U_0=0$, $U_1=1$, $U_{j+1}=c'U_j-U_{j-1}$, identifies
$|w|=U_j$ and $|t|\in\{U_{j-1},U_{j+1}\}$.
Since $|A-2|$ divides $w$, reduction modulo $c'-2$ or $c'+2$ gives
$j\ge c'-2$. For $j\ge3$, the retained coordinate satisfies
$|t|\ge U_{j-1}\ge(c'-1)^{j-2}$. Thus the following lower bound
holds for $B$ itself, without inserting $c'$ into its definition:
\[
 \log B\ge(c'-5)\log(c'-2)
          \ge(c-6)\log(c-3),\qquad c=|a|.
\]
No assertion that $a+1$ has the same sign as $a$ is needed; the estimate
uses only $c'\ge c-1$.

If $J_X\ne0$, the terms involving $h(\beta)$ and $N$ in the transfer
bound can be absorbed into $B^2/2$ once $B$ exceeds a fixed constant.
The bound with $k=1$ therefore gives
\[
                    \log B\le\frac75 h(\beta)+C_5.
\]
Increase the single $N$ threshold until
$(7/5)C_3 149^{-L_N}\le1$. The last upper bound becomes
$c+(c/3)\log c+C_5$, contradicting
$(c-6)\log(c-3)$ for every sufficiently large $c$.
Hence $\beta=n\in\tfrac12\Z$. For reduced $X=u/v$, $v>0$,
this implies $v^2\mid102$. Since $102$ is squarefree, $v=1$.
Thus soundness holds with one fixed ineffective threshold.

\section*{Completeness}
Fix an integer input and a parameter $N$ beyond the unchanged uniform
scalar-completeness threshold. In the scalar preparation choose the
positive scalar index $a$ sufficiently large and divisible by $6$, in
addition to its existing finite-place divisibilities and real recurrence
conditions. The source preparation permits arbitrary extra fixed
integer divisibility: the pinned upstream \texttt{uniform-existence.md},
P13--P14, uses compact-group recurrence in finite-index progressions and
chooses $a$ arbitrarily large and deeper at the finite support, with
$[a]P$ in a fixed neighborhood of $O$. It explicitly imposes no residue
condition $a=1$. Enlarging this modulus does not change $N$.
Set $A=a+1$ and $j=a-1=A-2$. For the sequence
$U_0=0$, $U_1=1$, $U_{r+1}=AU_r-U_{r-1}$ put
\[
                    b=U_j/(a-1),\qquad t=U_{j+1}.
\]
Modulo $a-1=A-2$, the recurrence gives $U_j\equiv j=0$; hence $b$ is
an integer. Because $a$ is divisible by $6$, $A$ and $j$ are odd and
$j\equiv2\pmod3$. Modulo $2$ the recurrence has period $0,1,1$,
so $U_j$ is odd. Therefore $b$ is positive and odd; by choosing $a$
large, $b>1$. The norm-one identity gives $\Delta=0$.
Take
\[
 Q_1=[\beta]P,\qquad Q_2=[a]P,\qquad Q_3=[t]R_0,
 \qquad Q_4=[(b-1)/2]R_0.
\]
The fourth point is a nonzero integer multiple of $R_0$, and
$4\nu(Q_4)+1=b$. All four points, the weight $R$, and the positive
trace $T$ are fixed before comparison-prime selection.

For the reference integer $h$ the new fourth numerator points are
\[
 U_4=[h(b-1)/2]R_0,\qquad V_4=[(h-1)(b-1)/2]R_0.
\]
The source Gaussian reciprocity lemma annihilates $R_0$ at the same
conjugate reduction-order moduli. Thus the unchanged $h\equiv0$ and
$h\equiv1$ congruences make this pair formal at its designated places.
The affine ratios there are $b$ to the required precision. All other
comparison and input residues are unchanged. Since the labels are now
fixed, the same compact-group recurrence supplies joint real neighborhoods
avoiding the finite chart exceptions, and the same fixed-data prime
premise applies to the rational coefficient list
$\Theta\cup(\Theta-T)$. Repetitions are merged with their incidences.

The factor exponents, contact squareclasses, and actual pole inputs are
unchanged in form. The source argument therefore proves their square
reduced denominators, prepared finite-place depths and squareclasses,
and real interval conditions. Its pole-realization premise supplies
witnesses. Filling the reciprocal pairs and affine ratios satisfies the
new quadratic equations. No input-dependent lower bound on $N$ has been
introduced. Together with soundness this proves the claimed equality of
input projections for every sufficiently large fixed $N$.

\section*{Algebraic count and scope}
The fourth root block uses seven field coordinates and nine field
quadrics. Its replacement uses five field coordinates and nine field
quadrics, reusing the same square and ratio outputs. The two Lucas
multiplication outputs and the aggregate retain their counts.
Consequently the proposed table has $204-4=200$ total rational variables,
with the same 186 rational quadratic residuals. This paragraph is an
algebraic specification, not a substitute for a complete table and its
independent coefficient check. The inherited hypotheses and ineffective
sufficient parameter threshold remain explicit.
\end{document}
